% =====================================================================
%  Appendix A: Single-Variable Calculus Quick Reference
% =====================================================================

\chapter{Single-Variable Calculus Quick Reference}
\label{app:single-variable-review}

This appendix is a reference for the single-variable facts used throughout the book.
It is meant to be scanned, not read like a chapter.

\[
\begin{array}{c|c}
\text{Single-variable idea} & \text{Multivariable version} \\ \hline
\text{derivative} & \text{partial derivative, directional derivative, differential}\\[1mm]
\text{linear approximation} & \text{tangent plane, Jacobian matrix}\\[1mm]
\text{substitution} & \text{change of variables, pullback}\\[1mm]
\text{definite integral} & \text{double, triple, line, and surface integrals}\\[1mm]
\text{Fundamental Theorem of Calculus} & \text{Green, Stokes, divergence, generalized Stokes}
\end{array}
\]

% ---------------------------------------------------------------------
\section{Limits and continuity}
\label{app:A.1}

\subsection*{Limits}

\[
\boxed{
    \lim_{x\to a}f(x)=L
}
\]
means that \(f(x)\) gets close to \(L\) as \(x\) gets close to \(a\), with \(x\ne a\).

The value \(f(a)\) does not have to exist.

\[
    \lim_{x\to a}f(x)
\]
exists exactly when the two one-sided limits agree:
\[
\boxed{
    \lim_{x\to a^-}f(x)=\lim_{x\to a^+}f(x)=L.
}
\]

\subsection*{Common limit laws}

If the limits exist, then
\[
\lim_{x\to a}(f+g)=\lim_{x\to a}f+\lim_{x\to a}g,
\]
\[
\lim_{x\to a}(fg)=
\left(\lim_{x\to a}f\right)
\left(\lim_{x\to a}g\right),
\]
and, if \(\lim_{x\to a}g(x)\ne0\),
\[
\lim_{x\to a}\frac{f}{g}
=
\frac{\lim_{x\to a}f}{\lim_{x\to a}g}.
\]

\subsection*{Useful standard limits}

\[
\boxed{
    \lim_{x\to0}\frac{\sin x}{x}=1
}
\]

\[
\boxed{
    \lim_{x\to0}\frac{1-\cos x}{x}=0
}
\]

\[
\boxed{
    \lim_{x\to0}\frac{e^x-1}{x}=1
}
\]

\[
\boxed{
    \lim_{x\to0}\frac{\ln(1+x)}{x}=1
}
\]

\subsection*{Continuity}

A function \(f\) is continuous at \(a\) if
\[
\boxed{
    \lim_{x\to a}f(x)=f(a).
}
\]

This requires three things:
\[
    f(a)\text{ is defined},
\]
\[
    \lim_{x\to a}f(x)\text{ exists},
\]
and
\[
    \lim_{x\to a}f(x)=f(a).
\]

\subsection*{Common continuous functions}

Polynomials are continuous everywhere.

Rational functions are continuous where their denominators are nonzero.

The functions
\[
    \sin x,\qquad \cos x,\qquad e^x
\]
are continuous everywhere.

The function
\[
    \ln x
\]
is continuous for
\[
    x>0.
\]

Compositions of continuous functions are continuous where they are defined.

\subsection*{Interval theorems}

\begin{theorem}[Intermediate Value Theorem]
If \(f\) is continuous on \([a,b]\), and \(N\) lies between \(f(a)\) and \(f(b)\), then
there is some \(c\in[a,b]\) such that
\[
    f(c)=N.
\]
\end{theorem}

\begin{theorem}[Extreme Value Theorem]
If \(f\) is continuous on a closed interval \([a,b]\), then \(f\) has an absolute maximum
and an absolute minimum on \([a,b]\).
\end{theorem}

\subsection*{Warnings}

A limit can exist even if \(f(a)\) is missing.

A function can be defined at \(a\) but still fail to be continuous there.

Continuity on an open interval does not guarantee an absolute maximum or minimum.

In several variables, checking only one path is not enough to prove a limit exists.

% ---------------------------------------------------------------------
\section{Derivatives and linear approximation}
\label{app:A.2}

\subsection*{Derivative}

\[
\boxed{
    f'(a)=\lim_{h\to0}\frac{f(a+h)-f(a)}{h}
}
\]
when the limit exists.

The derivative is:

\[
\begin{array}{c|c}
\text{Meaning} & \text{Description} \\ \hline
\text{slope} & \text{slope of the tangent line}\\
\text{rate} & \text{instantaneous rate of change}\\
\text{linear part} & \text{best first-order approximation near }a
\end{array}
\]

\subsection*{Basic derivative formulas}

\[
\frac{d}{dx}c=0
\]

\[
\frac{d}{dx}x^n=nx^{n-1}
\]

\[
\frac{d}{dx}e^x=e^x
\]

\[
\frac{d}{dx}\ln x=\frac1x
\qquad (x>0)
\]

\[
\frac{d}{dx}\sin x=\cos x
\]

\[
\frac{d}{dx}\cos x=-\sin x
\]

\[
\frac{d}{dx}\tan x=\sec^2x
\]

\subsection*{Derivative rules}

\[
\boxed{
    (f+g)'=f'+g'
}
\]

\[
\boxed{
    (cf)'=cf'
}
\]

\[
\boxed{
    (fg)'=f'g+fg'
}
\]

\[
\boxed{
    \left(\frac{f}{g}\right)'
    =
    \frac{f'g-fg'}{g^2}
    \qquad (g\ne0)
}
\]

\[
\boxed{
    \frac{d}{dx}f(g(x))=f'(g(x))g'(x)
}
\]
This is the chain rule.

\subsection*{Linear approximation}

If \(f\) is differentiable at \(a\), then near \(a\),
\[
\boxed{
    f(x)\approx f(a)+f'(a)(x-a).
}
\]

The tangent line is
\[
\boxed{
    L(x)=f(a)+f'(a)(x-a).
}
\]

\subsection*{Differential notation}

\[
\boxed{
    df=f'(x)\,dx.
}
\]

For a small input change \(dx\), the predicted output change is
\[
    df.
\]

Usually,
\[
    \Delta f\ne df,
\]
but for small changes,
\[
    \Delta f\approx df.
\]

\subsection*{Second derivative}

\[
    f''(x)=\frac{d}{dx}f'(x).
\]

\[
\begin{array}{c|c}
\text{Condition} & \text{Meaning} \\ \hline
f''(x)>0 & \text{graph is concave up}\\
f''(x)<0 & \text{graph is concave down}\\
f''(x)=0 & \text{possible inflection point}
\end{array}
\]

\subsection*{Critical points}

A critical point occurs where
\[
    f'(x)=0
\]
or where \(f'\) does not exist.

For absolute extrema on \([a,b]\), check:
\[
    \text{critical points in }(a,b),
    \qquad
    a,
    \qquad
    b.
\]

\subsection*{Warnings}

Differentiability implies continuity:
\[
\boxed{
    f\text{ differentiable at }a
    \Longrightarrow
    f\text{ continuous at }a.
}
\]

The converse is false.  For example,
\[
    f(x)=|x|
\]
is continuous at \(0\), but not differentiable at \(0\).

The derivative may fail because of:

\[
\begin{array}{c|c}
\text{Feature} & \text{Example} \\ \hline
\text{corner} & |x|\\
\text{vertical tangent} & \sqrt[3]{x}\text{ at }0\\
\text{jump} & \text{step function}\\
\text{oscillation} & x\sin(1/x)\text{ near }0
\end{array}
\]

% ---------------------------------------------------------------------
\section{The Fundamental Theorem of Calculus}
\label{app:A.3}

\subsection*{Definite integral}

\[
\boxed{
    \int_a^b f(x)\,dx
}
\]
is the signed accumulation of \(f\) from \(a\) to \(b\).

If \(f(x)\ge0\), it is area under the graph.

If \(f\) is a rate, then
\[
    \int_a^b f(x)\,dx
\]
is total change.

\subsection*{Basic properties}

\[
\int_a^b (f+g)\,dx
=
\int_a^b f\,dx+\int_a^b g\,dx
\]

\[
\int_a^b cf\,dx
=
c\int_a^b f\,dx
\]

\[
\int_a^b f\,dx
=
-\int_b^a f\,dx
\]

\[
\int_a^a f\,dx=0
\]

\[
\int_a^b f\,dx
=
\int_a^c f\,dx+\int_c^b f\,dx
\]

If
\[
    f(x)\ge0
\]
on \([a,b]\), then
\[
    \int_a^b f(x)\,dx\ge0.
\]

If
\[
    f(x)\le g(x)
\]
on \([a,b]\), then
\[
    \int_a^b f(x)\,dx\le \int_a^b g(x)\,dx.
\]

\subsection*{Fundamental Theorem of Calculus, Part I}

If \(f\) is continuous and
\[
    A(x)=\int_a^x f(t)\,dt,
\]
then
\[
\boxed{
    A'(x)=f(x).
}
\]

Accumulation differentiates back to the integrand.

\subsection*{Fundamental Theorem of Calculus, Part II}

If
\[
    F'(x)=f(x),
\]
then
\[
\boxed{
    \int_a^b f(x)\,dx=F(b)-F(a).
}
\]

Equivalently,
\[
\boxed{
    \int_a^b F'(x)\,dx=F(b)-F(a).
}
\]

This is the net change theorem.

\subsection*{Boundary form}

The oriented boundary of an interval is
\[
\boxed{
    \partial[a,b]=\{b\}-\{a\}.
}
\]

So the Fundamental Theorem can be read as
\[
\boxed{
    \int_{[a,b]} dF=\int_{\partial[a,b]}F.
}
\]

This is the one-dimensional model for the generalized Stokes theorem.

\subsection*{Average value}

The average value of \(f\) on \([a,b]\) is
\[
\boxed{
    f_{\operatorname{avg}}
    =
    \frac{1}{b-a}\int_a^b f(x)\,dx.
}
\]

\subsection*{Common antiderivatives}

\[
\int x^n\,dx=\frac{x^{n+1}}{n+1}+C
\qquad (n\ne -1)
\]

\[
\int \frac1x\,dx=\ln|x|+C
\]

\[
\int e^x\,dx=e^x+C
\]

\[
\int \sin x\,dx=-\cos x+C
\]

\[
\int \cos x\,dx=\sin x+C
\]

\[
\int \sec^2x\,dx=\tan x+C
\]

\subsection*{Warnings}

Do not forget the constant \(C\) in an indefinite integral.

Do not put \(+C\) on a definite integral.

The definite integral
\[
    \int_a^b f(x)\,dx
\]
depends on the orientation of the interval.  Reversing the endpoints changes the sign.

% ---------------------------------------------------------------------
\section{Substitution and integration by parts}
\label{app:A.4}

\subsection*{Substitution}

Substitution reverses the chain rule.

If
\[
    u=g(x),
\]
then
\[
\boxed{
    du=g'(x)\,dx.
}
\]

\[
\boxed{
    \int f(g(x))g'(x)\,dx=\int f(u)\,du.
}
\]

For definite integrals,
\[
\boxed{
    \int_a^b f(g(x))g'(x)\,dx
    =
    \int_{g(a)}^{g(b)} f(u)\,du.
}
\]

\subsection*{Substitution checklist}

\[
\begin{array}{ll}
1. & \text{Choose }u=g(x).\\
2. & \text{Compute }du=g'(x)\,dx.\\
3. & \text{Replace the integrand and differential.}\\
4. & \text{For definite integrals, change the limits.}\\
5. & \text{For indefinite integrals, substitute back to }x.
\end{array}
\]

\subsection*{Common substitution patterns}

\[
\int f(ax+b)\,dx
\quad\text{use}\quad
u=ax+b.
\]

\[
\int f(x^2)\,2x\,dx
\quad\text{use}\quad
u=x^2.
\]

\[
\int f(\sin x)\cos x\,dx
\quad\text{use}\quad
u=\sin x.
\]

\[
\int f(\cos x)(-\sin x)\,dx
\quad\text{use}\quad
u=\cos x.
\]

\subsection*{Warning}

The substitution
\[
    u=g(x)
\]
requires the matching factor
\[
    g'(x)\,dx.
\]

For example,
\[
    \int 2x e^{x^2}\,dx
\]
is directly handled by \(u=x^2\), but
\[
    \int e^{x^2}\,dx
\]
is not.

\subsection*{Integration by parts}

Integration by parts reverses the product rule.

\[
\boxed{
    \int u\,dv=uv-\int v\,du.
}
\]

For definite integrals,
\[
\boxed{
    \int_a^b u\,dv
    =
    \left[uv\right]_a^b-\int_a^b v\,du.
}
\]

\subsection*{Useful choices}

For
\[
    \int x e^x\,dx,
\]
take
\[
    u=x,
    \qquad
    dv=e^x\,dx.
\]

For
\[
    \int x\cos x\,dx,
\]
take
\[
    u=x,
    \qquad
    dv=\cos x\,dx.
\]

For
\[
    \int \ln x\,dx,
\]
take
\[
    u=\ln x,
    \qquad
    dv=dx.
\]

\subsection*{Warning}

Integration by parts creates a boundary term:
\[
    [uv]_a^b.
\]
Do not drop it.

% ---------------------------------------------------------------------
\section{Parametric and polar curves}
\label{app:A.5}

\subsection*{Parametric curves}

A plane curve can be written as
\[
\boxed{
    \mathbf r(t)=(x(t),y(t)).
}
\]

A space curve can be written as
\[
\boxed{
    \mathbf r(t)=(x(t),y(t),z(t)).
}
\]

The derivative is
\[
\boxed{
    \mathbf r'(t)=(x'(t),y'(t))
}
\]
in the plane, and
\[
\boxed{
    \mathbf r'(t)=(x'(t),y'(t),z'(t))
}
\]
in space.

\subsection*{Velocity, speed, and acceleration}

\[
\boxed{
    \mathbf v(t)=\mathbf r'(t)
}
\]

\[
\boxed{
    \text{speed}=|\mathbf r'(t)|
}
\]

\[
\boxed{
    \mathbf a(t)=\mathbf r''(t)
}
\]

In the plane,
\[
\boxed{
    |\mathbf r'(t)|
    =
    \sqrt{(x'(t))^2+(y'(t))^2}.
}
\]

In space,
\[
\boxed{
    |\mathbf r'(t)|
    =
    \sqrt{(x'(t))^2+(y'(t))^2+(z'(t))^2}.
}
\]

\subsection*{Arc length}

For a curve \(\mathbf r(t)\), \(a\le t\le b\),
\[
\boxed{
    L=\int_a^b |\mathbf r'(t)|\,dt.
}
\]

\subsection*{Slope of a parametric plane curve}

If
\[
    x'(t)\ne0,
\]
then
\[
\boxed{
    \frac{dy}{dx}
    =
    \frac{dy/dt}{dx/dt}.
}
\]

If
\[
    x'(t)=0
    \quad\text{and}\quad
    y'(t)\ne0,
\]
the tangent is vertical.

\subsection*{Common parametrizations}

Line segment from \(P\) to \(Q\):
\[
\boxed{
    \mathbf r(t)=P+t(Q-P),
    \qquad
    0\le t\le1.
}
\]

Circle of radius \(R\), counterclockwise:
\[
\boxed{
    \mathbf r(t)=(R\cos t,R\sin t),
    \qquad
    0\le t\le2\pi.
}
\]

Circle of radius \(R\), clockwise:
\[
\boxed{
    \mathbf r(t)=(R\cos t,-R\sin t),
    \qquad
    0\le t\le2\pi.
}
\]

Helix:
\[
\boxed{
    \mathbf r(t)=(R\cos t,R\sin t,ct).
}
\]

\subsection*{Polar coordinates}

\[
\boxed{
    x=r\cos\theta,
    \qquad
    y=r\sin\theta.
}
\]

\[
\boxed{
    r^2=x^2+y^2.
}
\]

\[
\boxed{
    \tan\theta=\frac{y}{x}
}
\]
with quadrant care.

A polar curve
\[
    r=f(\theta)
\]
can be parametrized by
\[
\boxed{
    \mathbf r(\theta)
    =
    \bigl(f(\theta)\cos\theta,\ f(\theta)\sin\theta\bigr).
}
\]

\subsection*{Polar area}

If \(r=f(\theta)\) traces the region once for
\[
    \alpha\le\theta\le\beta,
\]
then
\[
\boxed{
    A=\frac12\int_\alpha^\beta f(\theta)^2\,d\theta.
}
\]

\subsection*{Polar arc length}

For \(r=f(\theta)\),
\[
\boxed{
    L=
    \int_\alpha^\beta
    \sqrt{
        r^2+\left(\frac{dr}{d\theta}\right)^2
    }
    \,d\theta.
}
\]

\subsection*{Warning}

A parametrization remembers orientation and speed.  The same geometric curve may give
different line integrals if it is traced in the opposite direction or traced multiple times.

% ---------------------------------------------------------------------
\section{Infinite series and Taylor polynomials}
\label{app:A.6}

\subsection*{Sequences}

A sequence is a list
\[
    a_1,a_2,a_3,\ldots.
\]

It converges to \(L\) if
\[
\boxed{
    \lim_{n\to\infty}a_n=L.
}
\]

\subsection*{Series}

A series is an infinite sum:
\[
\boxed{
    \sum_{n=1}^{\infty}a_n.
}
\]

Its partial sums are
\[
    s_N=a_1+a_2+\cdots+a_N.
\]

The series converges if
\[
    \lim_{N\to\infty}s_N
\]
exists.

\subsection*{Geometric series}

\[
\boxed{
    \sum_{n=0}^{\infty} ar^n=\frac{a}{1-r}
    \qquad
    (|r|<1).
}
\]

If
\[
    |r|\ge1,
\]
the geometric series diverges.

\subsection*{Basic convergence facts}

If
\[
    \sum a_n
\]
converges, then
\[
\boxed{
    a_n\to0.
}
\]

The converse is false.  The harmonic series
\[
    \sum_{n=1}^{\infty}\frac1n
\]
diverges, even though
\[
    \frac1n\to0.
\]

\subsection*{\(p\)-series}

\[
\boxed{
    \sum_{n=1}^{\infty}\frac1{n^p}
    \text{ converges if }p>1
}
\]

\[
\boxed{
    \sum_{n=1}^{\infty}\frac1{n^p}
    \text{ diverges if }p\le1.
}
\]

\subsection*{Power series}

A power series centered at \(a\) has the form
\[
\boxed{
    \sum_{n=0}^{\infty}c_n(x-a)^n.
}
\]

A power series usually converges only on an interval.

\subsection*{Common Taylor series}

\[
\boxed{
    e^x
    =
    1+x+\frac{x^2}{2!}+\frac{x^3}{3!}+\cdots
}
\]

\[
\boxed{
    \sin x
    =
    x-\frac{x^3}{3!}+\frac{x^5}{5!}-\cdots
}
\]

\[
\boxed{
    \cos x
    =
    1-\frac{x^2}{2!}+\frac{x^4}{4!}-\cdots
}
\]

\[
\boxed{
    \frac1{1-x}
    =
    1+x+x^2+x^3+\cdots
    \qquad
    (|x|<1)
}
\]

\[
\boxed{
    \ln(1+x)
    =
    x-\frac{x^2}{2}+\frac{x^3}{3}-\frac{x^4}{4}+\cdots
}
\]
valid for
\[
    -1<x\le1.
\]

\subsection*{Taylor polynomial}

The degree \(n\) Taylor polynomial of \(f\) at \(a\) is
\[
\boxed{
    T_n(x)
    =
    f(a)
    +
    f'(a)(x-a)
    +
    \frac{f''(a)}{2!}(x-a)^2
    +
    \cdots
    +
    \frac{f^{(n)}(a)}{n!}(x-a)^n.
}
\]

Degree \(1\):
\[
\boxed{
    T_1(x)=f(a)+f'(a)(x-a).
}
\]
This is the linear approximation.

Degree \(2\):
\[
\boxed{
    T_2(x)=f(a)+f'(a)(x-a)+\frac12 f''(a)(x-a)^2.
}
\]
This includes curvature.

\subsection*{Taylor's theorem with remainder}

If \(f\) has \(n+1\) derivatives near \(a\), then
\[
\boxed{
    f(x)=T_n(x)+R_n(x).
}
\]

One common form of the remainder is
\[
\boxed{
    R_n(x)
    =
    \frac{f^{(n+1)}(c)}{(n+1)!}(x-a)^{n+1}
}
\]
for some \(c\) between \(a\) and \(x\).

\subsection*{Warning}

A Taylor polynomial is always a local approximation.

A Taylor series may fail to equal the original function, even when the function has
derivatives of all orders.

\subsection*{Multivariable connection}

The one-variable approximation
\[
    f(a+h)\approx f(a)+f'(a)h
\]
becomes
\[
    f(\mathbf p+\mathbf h)
    \approx
    f(\mathbf p)+Df_{\mathbf p}(\mathbf h).
\]

The one-variable quadratic approximation
\[
    f(a+h)\approx f(a)+f'(a)h+\frac12 f''(a)h^2
\]
becomes
\[
    f(\mathbf p+\mathbf h)
    \approx
    f(\mathbf p)
    +
    Df_{\mathbf p}(\mathbf h)
    +
    \frac12\mathbf h^T H_f(\mathbf p)\mathbf h.
\]

% ---------------------------------------------------------------------
\section*{One-page formula summary}

\[
\boxed{
    f'(a)=\lim_{h\to0}\frac{f(a+h)-f(a)}{h}
}
\]

\[
\boxed{
    L(x)=f(a)+f'(a)(x-a)
}
\]

\[
\boxed{
    df=f'(x)\,dx
}
\]

\[
\boxed{
    \int_a^b F'(x)\,dx=F(b)-F(a)
}
\]

\[
\boxed{
    \int f(g(x))g'(x)\,dx=\int f(u)\,du
}
\]

\[
\boxed{
    \int u\,dv=uv-\int v\,du
}
\]

\[
\boxed{
    L_{\text{curve}}=\int_a^b|\mathbf r'(t)|\,dt
}
\]

\[
\boxed{
    A_{\text{polar}}=\frac12\int_\alpha^\beta r^2\,d\theta
}
\]

\[
\boxed{
    \sum_{n=0}^{\infty}ar^n=\frac{a}{1-r}
    \quad (|r|<1)
}
\]

\[
\boxed{
    T_n(x)=\sum_{k=0}^{n}\frac{f^{(k)}(a)}{k!}(x-a)^k
}
\]

\[
\boxed{
    \partial[a,b]=\{b\}-\{a\}
}
\]

\[
\boxed{
    \int_{[a,b]}dF=\int_{\partial[a,b]}F
}
\]